\documentclass[conference]{IEEEtran}

\usepackage[T1]{fontenc}
\usepackage[utf8]{inputenc}
\usepackage{cite}
\usepackage{amsmath,amssymb}
\usepackage{booktabs}
\usepackage{tabularx}
\usepackage{array}
\usepackage{enumitem}
\usepackage{hyperref}
\usepackage{xcolor}

\hypersetup{
  colorlinks=true,
  linkcolor=black,
  citecolor=black,
  urlcolor=blue
}

\title{A CAPIF-Aware AI-Native Telco Orchestration Platform for\\Intent-to-API Planning Across CAMARA and NEF Services}

\author{
\IEEEauthorblockN{T\"urkan Do\u{g}a Durak}
\IEEEauthorblockA{
CTTC Internship Project\\
Barcelona, Spain\\
\texttt{turkan.doga.durak@example.com}
}
}

\begin{document}

\maketitle

\begin{abstract}
This paper presents the design and implementation of a mini telco orchestration platform that combines Common API Framework (CAPIF) integration, CAMARA-aligned service semantics, NEF-backed telecom realization, and large language model (LLM)-assisted intent planning. The system is implemented as a FastAPI-based control plane with a product-style operator interface, embedded OpenAPI console, live CAPIF onboarding and discovery, provider publication, persistent NEF data services, and guarded orchestration execution. The platform supports both deterministic and LLM-assisted planning modes and validates generated API payloads before execution. A live Gemini-backed orchestration test confirms that natural-language service requests can be translated into structured QoD execution plans under runtime grounding and deterministic validation. The resulting system provides a practical baseline for reliability-aware intent-to-API orchestration research in 5G and beyond-5G service exposure environments.
\end{abstract}

\begin{IEEEkeywords}
CAPIF, CAMARA, NEF, 5G service exposure, intent-based orchestration, LLM planning, telecom APIs, OpenCAPIF
\end{IEEEkeywords}

\section{Introduction}

Programmable network exposure is becoming a central design principle in modern telecom systems. CAPIF provides a standardized governance and discovery plane for exposed APIs, CAMARA defines developer-facing service semantics, and NEF-based providers bridge those interfaces to telecom data and network capabilities. At the same time, large language models introduce a new opportunity: translating high-level intents into structured API workflows without requiring the operator to manually compose payloads and endpoint sequences.

This project implements a working telco orchestration platform at that intersection. Rather than remaining at the conceptual architecture level, the system integrates live OpenCAPIF onboarding and discovery, live provider publication, live QoD and Location Retrieval execution, an embedded operator interface, and an LLM-assisted planning layer protected by deterministic validation.

The platform was developed not as an isolated mockup, but as a usable product-style control surface for demonstrating and evaluating end-to-end telecom orchestration behavior. This makes it suitable both for technical demonstrations and for thesis-oriented experimentation on intent-to-API translation reliability.

\section{System Objectives and Contributions}

The platform was built with four primary objectives:

\begin{itemize}[leftmargin=1.1em]
\item to unify CAPIF onboarding, discovery, and publication inside one operator-facing control plane,
\item to expose CAMARA-aligned telecom capabilities through a live orchestration backend,
\item to support both deterministic and LLM-assisted planning for natural-language requests,
\item to ensure that execution remains guarded by deterministic validation and schema enforcement.
\end{itemize}

Its main contributions are:

\begin{itemize}[leftmargin=1.1em]
\item a working CAPIF-aware telco orchestration backend,
\item a pluggable LLM planning layer grounded by local service metadata and CAPIF discovery,
\item a live provider integration across QoD, Location Retrieval, and NEF Monitoring Event,
\item a persistent and automatically seeded NEF data layer,
\item a product-style front end with an embedded API console for live demonstrations.
\end{itemize}

\section{Platform Architecture}

The implemented system follows a layered control-plane architecture composed of user interaction, orchestration, planning, exposure governance, provider execution, and persistence layers.

\subsection{High-Level Components}

\begin{table}[h]
\caption{Implemented Platform Components}
\label{tab:components}
\centering
\small
\begin{tabularx}{\linewidth}{>{\raggedright\arraybackslash}p{0.24\linewidth} X}
\toprule
\textbf{Component} & \textbf{Responsibility} \\
\midrule
Operator Interface & Product-style landing page with embedded OpenAPI console and orchestration access \\
FastAPI Backend & Route handling, orchestration control, validation, and execution dispatch \\
Intent Engine & Deterministic parsing of user requests into structured execution steps \\
LLM Planner & Gemini-based structured planning under catalog and CAPIF grounding \\
CAPIF Client & Invoker onboarding, service discovery, and provider publication \\
Provider Layer & QoD, Location Retrieval, and NEF Monitoring Event service realization \\
Persistence Layer & MongoDB-backed NEF demo data with automatic initialization and persistence \\
\bottomrule
\end{tabularx}
\end{table}

\subsection{Repository Structure}

\begin{table}[h]
\caption{Key Repository Paths}
\label{tab:repo}
\centering
\small
\begin{tabularx}{\linewidth}{>{\ttfamily\raggedright\arraybackslash}p{0.33\linewidth} X}
\toprule
\textbf{Path} & \textbf{Role} \\
\midrule
backend/main.py & FastAPI entrypoint and primary route registration \\
backend/intent\_engine.py & Deterministic intent parsing and execution-step construction \\
backend/llm\_planner.py & LLM-assisted planning and grounding logic \\
backend/capif/client.py & CAPIF register, onboarding, discovery, and publish client \\
backend/models.py & Request, response, and payload schemas \\
frontend/index.html & Product-style operator landing page with embedded API console \\
providers/NEF/ & NEF Monitoring Event stack and Mongo persistence \\
providers/CamaraLocationRetrieval/ & CAMARA Location Retrieval service \\
providers/camara-QoD/ & CAMARA Quality-on-Demand service \\
docker-compose.capif.yml & Main deployment attached to the OpenCAPIF network \\
\bottomrule
\end{tabularx}
\end{table}

\section{Backend Orchestration Design}

\subsection{Core API Surface}

The backend exposes the control and orchestration endpoints listed in Table~\ref{tab:endpoints}.

\begin{table}[h]
\caption{Main Backend Endpoints}
\label{tab:endpoints}
\centering
\small
\begin{tabularx}{\linewidth}{>{\raggedright\arraybackslash}p{0.12\linewidth} >{\ttfamily\raggedright\arraybackslash}p{0.27\linewidth} X}
\toprule
\textbf{Method} & \textbf{Path} & \textbf{Purpose} \\
\midrule
GET & /health & Runtime health and status summary \\
GET & /catalog/services & Local CAMARA-aligned capability catalog \\
POST & /orchestrate & Intent-to-execution planning endpoint \\
GET & /capif/status & CAPIF and register connectivity check \\
GET & /capif/invoker & Stored invoker identity details \\
POST & /capif/invokers/onboard & Live invoker onboarding against CAPIF \\
GET & /capif/discover & Live service discovery through CAPIF \\
POST & /capif/publish & Provider publication endpoint \\
POST & /capif/publish/catalog & Publication of catalog services into CAPIF \\
\bottomrule
\end{tabularx}
\end{table}

\subsection{Supported Telecom Capabilities}

The current platform scope includes four orchestration targets:

\begin{itemize}[leftmargin=1.1em]
\item \textbf{Quality on Demand (QoD)} for temporary QoS session creation,
\item \textbf{QoS Profiles} for profile inspection,
\item \textbf{QoS Provisioning} for long-lived QoS assignment,
\item \textbf{Location Retrieval} for device location access through NEF-backed monitoring.
\end{itemize}

\subsection{Planning Modes}

The \texttt{POST /orchestrate} endpoint supports three planning modes:

\begin{itemize}[leftmargin=1.1em]
\item \texttt{deterministic}: always uses the rule-based planner,
\item \texttt{llm-assisted}: requires the LLM planner and fails explicitly if planning fails,
\item \texttt{auto}: attempts the LLM planner first and falls back to deterministic planning if needed.
\end{itemize}

This design keeps the AI layer optional yet operational, which is particularly useful for both live demos and controlled experimentation.

\section{CAPIF Integration Workflow}

\subsection{CAPIF Roles in the Platform}

CAPIF provides the exposure governance plane used by the system. The key roles are:

\begin{itemize}[leftmargin=1.1em]
\item \textbf{Invoker}: the mini orchestration platform or provider client consuming exposed APIs,
\item \textbf{Provider / AEF}: the service implementation exposing an API through CAPIF,
\item \textbf{CAPIF Core}: the OpenCAPIF services managing discovery and onboarding,
\item \textbf{Register}: the bootstrap service used for identity and credential preparation.
\end{itemize}

\subsection{Invoker Onboarding Flow}

The implemented onboarding flow in \texttt{backend/capif/client.py} performs the following steps:

\begin{enumerate}[leftmargin=1.1em]
\item authenticate against the register service,
\item ensure a demo user exists,
\item obtain credential material for that user,
\item generate a local private key and certificate signing request,
\item submit the onboarding request to CAPIF Core,
\item receive an \texttt{apiInvokerId} and invoker certificate,
\item store runtime invoker details for later discovery and execution.
\end{enumerate}

This workflow enables the platform to operate as a real CAPIF invoker rather than as a static mock client.

\subsection{Certificate and Credential Exchange}

The CAPIF onboarding sequence is also the certificate-exchange sequence of the platform. In operational terms, the platform does not begin with a long-lived client certificate already provisioned. Instead, it creates one dynamically during onboarding.

The implemented credential path is:

\begin{enumerate}[leftmargin=1.1em]
\item the platform authenticates against the register service,
\item the register service returns bootstrap authentication material,
\item the platform generates an RSA private key locally,
\item the platform generates a certificate signing request (CSR),
\item the CSR is submitted in the CAPIF onboarding request,
\item CAPIF returns an invoker certificate chain bound to the new invoker identity,
\item the platform stores the private key, certificate, and trust material under the runtime CAPIF directory,
\item later CAPIF requests can use this runtime certificate material for authenticated interactions.
\end{enumerate}

This is important for the thesis and the presentation because it shows that CAPIF is not only a catalog or registry service. It is also the trust-establishment layer that enables controlled access to exposed telecom APIs.

\subsection{Discovery Flow}

The discovery workflow is implemented by \texttt{GET /capif/discover}:

\begin{enumerate}[leftmargin=1.1em]
\item read the stored invoker identity,
\item extract the invoker identifier,
\item query the CAPIF discovery endpoint,
\item return live service descriptions from the registry.
\end{enumerate}

In the current platform deployment, discovery returns published services including QoD, Location Retrieval, NEF Monitoring Event, and the default OpenCAPIF demo service.

\subsection{Provider Publication Flow}

Provider publication follows the standard sequence:

\begin{enumerate}[leftmargin=1.1em]
\item start the provider service,
\item prepare the OpenAPI descriptor,
\item onboard the provider identity if needed,
\item publish API metadata into CAPIF,
\item verify that the API appears in CAPIF discovery results.
\end{enumerate}

Conceptually, CAPIF therefore contributes three distinct functions to the platform:

\begin{itemize}[leftmargin=1.1em]
\item \textbf{identity onboarding}, because invokers and providers need runtime credentials,
\item \textbf{service governance}, because APIs become visible only after publication,
\item \textbf{service discovery}, because the orchestration platform can reason about what is currently exposed.
\end{itemize}

\section{CAMARA and NEF Service Realization}

\subsection{Role of CAMARA}

CAMARA defines the exposed service semantics consumed by the orchestration platform. In this system:

\begin{itemize}[leftmargin=1.1em]
\item CAPIF governs onboarding and discovery,
\item CAMARA defines the API-facing service contracts,
\item provider services bridge these APIs to telecom behavior and data.
\end{itemize}

In this project, CAMARA is the language of service exposure. It defines what a QoD request should look like, what a location retrieval request should return, and what the application developer or platform operator sees at the API layer. This makes CAMARA the semantic layer of the system, while CAPIF remains the governance and exposure layer.

\subsection{How CAMARA Was Connected to the Platform}

The current platform integrates CAMARA-aligned services through FRONT research group provider implementations. The integration path is:

\begin{enumerate}[leftmargin=1.1em]
\item select the target CAMARA service family, such as QoD or Location Retrieval,
\item start the corresponding provider service,
\item align the backend catalog with the service identifier, operations, and payload structure,
\item publish or expose the provider through CAPIF-aware integration,
\item allow the orchestration backend to select the service from intent analysis,
\item validate and forward the request to the provider-facing execution path.
\end{enumerate}

As a result, the orchestration backend does not invent telecom semantics on its own. It reasons over a constrained set of known service contracts that are already aligned with CAMARA or provider OpenAPI descriptions.

\subsection{Role of NEF}

The NEF Monitoring Event provider acts as the telecom-facing realization layer for location workflows. It:

\begin{itemize}[leftmargin=1.1em]
\item receives a monitoring request,
\item maps phone numbers to IMSI identifiers,
\item reads location information from MongoDB,
\item enriches responses with polygonal area information,
\item returns a 3GPP-style monitoring result that upstream services can transform.
\end{itemize}

In simpler terms, NEF is the point where abstract API exposure meets telecom-side state. CAPIF and CAMARA make the service visible and consumable, but NEF is the component that actually connects the request to network-related information such as monitoring state, identifier mapping, and location data.

\subsection{Location Retrieval End-to-End Flow}

The live Location Retrieval path executes as follows:

\begin{enumerate}[leftmargin=1.1em]
\item the client calls the Location Retrieval provider,
\item the provider constructs a 3GPP Monitoring Event subscription,
\item the provider obtains CAPIF security context and access tokens,
\item the provider calls the NEF Monitoring Event API,
\item the NEF resolves identifiers and reads location plus polygon data,
\item the response is normalized into CAMARA Location Retrieval format,
\item the final location response is returned to the caller.
\end{enumerate}

\subsection{QoD Flow}

The QoD path is simpler:

\begin{enumerate}[leftmargin=1.1em]
\item the user intent is mapped to the \texttt{qod} capability,
\item the planner generates a session creation payload,
\item the QoD provider handles the session request,
\item the provider result is returned or validated in dry-run mode.
\end{enumerate}

\section{Persistence and Reliability Engineering}

\subsection{Persistent NEF Data}

During earlier integration stages, the NEF stack lost location and polygon records after recreation. This created operational instability for the Location Retrieval chain. The issue was resolved through:

\begin{itemize}[leftmargin=1.1em]
\item a named Docker volume mounted at \texttt{/data/db},
\item an automatic Mongo initialization script mounted at \texttt{/docker-entrypoint-initdb.d/mongo-init.js}.
\end{itemize}

The resulting behavior is:

\begin{enumerate}[leftmargin=1.1em]
\item the first deployment seeds the required data automatically,
\item standard restarts preserve the existing dataset,
\item the NEF-backed location pipeline survives stack recreation without manual reseeding.
\end{enumerate}

\subsection{Validation Guardrails}

Execution never depends solely on LLM output. All generated payloads are validated using backend-side schemas and structured execution models. This ensures:

\begin{itemize}[leftmargin=1.1em]
\item correct field names,
\item correct structural nesting,
\item presence of required values,
\item deterministic failure before provider invocation if a payload is malformed.
\end{itemize}

This separation between probabilistic planning and deterministic enforcement is a central architectural property of the platform.

\section{LLM-Assisted Planning Layer}

\subsection{Planning Strategy}

The file \texttt{backend/llm\_planner.py} implements the AI planning path. It follows this procedure:

\begin{enumerate}[leftmargin=1.1em]
\item load the local service catalog,
\item optionally incorporate live CAPIF discovery results,
\item build a constrained prompt containing allowed services, operations, and payload expectations,
\item call the configured model provider,
\item parse the returned JSON plan,
\item transform the output into an internal execution step,
\item pass the result into deterministic validation.
\end{enumerate}

The planner therefore proposes structured actions but never executes them directly.

\subsection{Live Gemini Orchestration Result}

After introducing a new Gemini project and API key, the LLM-assisted path was validated with a live end-to-end request using the intent:

\begin{quote}
Boost video QoS for +905551112233 for 15 minutes.
\end{quote}

The request was submitted to \texttt{POST /orchestrate} in \texttt{llm-assisted} mode with \texttt{dry\_run=true}. The backend returned HTTP 200 and produced the following execution semantics:

\begin{itemize}[leftmargin=1.1em]
\item selected service: \texttt{qod},
\item selected operation: \texttt{createSession},
\item target path: \texttt{/quality-on-demand/vwip/sessions},
\item device identifier: \texttt{+905551112233},
\item duration: \texttt{900} seconds,
\item QoS profile: \texttt{QOS\_E\_STREAMING}.
\end{itemize}

This confirms that the LLM planner now succeeds as a live runtime component rather than acting only as a fallback candidate under external quota failures.

\subsection{Productized Front End}

To support direct presentation and operation, the front end was also converted from a demo-style landing page into a product-style control surface. The revised interface:

\begin{itemize}[leftmargin=1.1em]
\item removes presentation-specific or informal wording,
\item presents the system as a telecom control plane,
\item embeds the live API console directly in the primary workspace,
\item provides a concise operator workflow for onboarding, discovery, and orchestration.
\end{itemize}

\section{Current End-to-End System Operation}

The running platform currently behaves as follows:

\begin{enumerate}[leftmargin=1.1em]
\item the operator enters through the product-style landing page,
\item the embedded API console exposes health, CAPIF, catalog, and orchestration routes,
\item the platform can onboard itself as a CAPIF invoker and obtain runtime certificate material,
\item the platform can inspect its local service catalog and attempt CAPIF discovery,
\item deterministic or LLM-assisted planning converts a natural-language request into a structured execution step,
\item the execution step is validated before any live provider call is made,
\item the validated request is then routed toward the relevant CAMARA or NEF-backed provider path.
\end{enumerate}

This means the system is not a single monolithic service. It is a connected control plane composed of governance, planning, validation, and provider execution layers.

\section{Operator Workflow and API Console}

The intended live presentation flow is:

\begin{enumerate}[leftmargin=1.1em]
\item open the landing page at \texttt{http://localhost:8000/},
\item verify backend health and CAPIF connectivity,
\item inspect the local service catalog,
\item perform invoker onboarding,
\item review discovered services through CAPIF,
\item submit orchestration requests through the embedded API console.
\end{enumerate}

Representative orchestration bodies include:

\begin{itemize}[leftmargin=1.1em]
\item \texttt{\{"intent": "Boost video QoS for +905551112233 for 15 minutes", "dry\_run": true, "orchestration\_mode": "llm-assisted"\}},
\item \texttt{\{"intent": "Retrieve the latest location for +905551112233", "dry\_run": true, "orchestration\_mode": "llm-assisted"\}},
\item \texttt{\{"intent": "Find the best available QoS profile for +905551112233", "dry\_run": true, "orchestration\_mode": "auto"\}}.
\end{itemize}

\section{Current Status and Next Steps}

At the current stage, the platform provides:

\begin{itemize}[leftmargin=1.1em]
\item live CAPIF invoker onboarding,
\item live CAPIF service discovery,
\item live provider publication visibility through CAPIF,
\item live QoD and NEF-backed Location Retrieval service paths,
\item persistent NEF data with automatic initialization,
\item a working Gemini-backed LLM planner,
\item a product-style operator interface with an embedded API console.
\end{itemize}

The next logical extensions include richer policy validation, additional CAMARA services, telemetry and evaluation instrumentation, and a formal reliability scoring layer for structured intent-to-API translation experiments.

\section{Conclusion}

This paper presented a working AI-native telco orchestration platform that integrates CAPIF governance, CAMARA service semantics, NEF-backed telecom realization, deterministic validation, and LLM-assisted planning into one deployable control plane. The system goes beyond mock orchestration by supporting live onboarding, discovery, provider publication, and execution-ready planning. The successful Gemini-backed orchestration result demonstrates that natural-language telecom requests can be translated into structured API actions under explicit runtime constraints. As a result, the platform forms a practical baseline for future work on trustworthy intent-driven service exposure in 5G and beyond-5G environments.

\end{document}
